\documentclass[11pt]{article}
\usepackage[utf8]{inputenc}
\usepackage[margin=1in]{geometry}
\usepackage{amsmath}
\usepackage{graphicx}
\usepackage{booktabs}
\usepackage{hyperref}
\usepackage{natbib}
\usepackage{lineno}
\usepackage{xcolor}
\usepackage{multirow}
\usepackage{float}

\linenumbers

\title{The Breath Shape Controls Intonation of Mouse Ultrasonic Vocalizations via the iRO Brainstem Circuit}
\author{
  Anonymous Authors\\
  \textit{Submitted to eLife}
}
\date{}

\begin{document}
\maketitle

\begin{abstract}
Mouse ultrasonic vocalizations (USVs) are widely used to study social communication and neurodevelopmental disorders, yet the mechanisms controlling their pitch patterns (intonation) remain poorly understood. It is generally assumed that dynamic changes in laryngeal tension are the sole mechanism for pitch modulation. Here we demonstrate that expiratory breath waveform modulation provides an additional, primary intonation mechanism for the majority of USV syllable types. Using simultaneous plethysmography and acoustic recordings in $n = 6$ male mice (1,850 vocalizations analyzed), we show that 6 of 10 syllable types exhibit positive intonation---pitch positively correlated with instantaneous expiratory airflow ($r = 0.52 \pm 0.08$ for Down FM; $p < 0.001$, one-way ANOVA with Sidak's correction). Diaphragm EMG recordings reveal that re-activation of the inspiratory diaphragm muscle during expiration creates the breath waveform modulations underlying positive intonation. Optogenetic activation of the intermediate reticular oscillator (iRO; Penk$^+$/Vglut2$^+$ neurons) is sufficient to entrain breathing and evoke USVs comprising most endogenous syllable types, producing primarily positive intonation patterns. USVs occurred during rapid sniff episodes (mean 7.5~Hz vs.\ 8.5--10~Hz for non-USV sniffs; $p = 0.03$, paired $t$-test), with 88\% of breaths containing a single syllable. We propose a two-mechanism model of intonation: positive (breath-driven, iRO-mediated) and negative (laryngeal tension-driven), which can combine within single vocalizations to generate the full diversity of the murine vocal repertoire.
\end{abstract}

\section{Introduction}

Mouse ultrasonic vocalizations (USVs) are a primary model system for studying mammalian social communication and its disruption in neurodevelopmental disorders \citep{holy2005ultrasonic, portfors2007types, hage2013murine}. Adult male mice produce rich vocal repertoires containing 10 or more distinct syllable types during courtship, each defined by characteristic pitch (frequency) contours spanning 40--120~kHz \citep{grimsley2011auditory, arriaga2012mouse}. The mechanisms generating this diversity of pitch patterns---the ``intonation'' of USVs---are central to understanding vocal motor control.

The prevailing model attributes pitch modulation entirely to dynamic changes in laryngeal muscle tension \citep{riede2011subglottal, riede2017mechanisms}. In this framework, the cricothyroid and thyroarytenoid muscles adjust vocal fold tension to modulate frequency, while expiratory airflow serves only to power phonation and control loudness \citep{titze1988regulation, alipour2002pressure}. However, a paradox exists: injecting air below the larynx increases pitch in rodents, yet natural breath airflow does not strongly predict pitch in naive analyses \citep{riede2011subglottal}. This discrepancy suggests that the relationship between airflow and pitch may be more nuanced than previously appreciated.

The intermediate reticular oscillator (iRO), a recently described brainstem central pattern generator, coordinates breathing and phonation \citep{delpouve2023intermediate, wei2022neural}. The iRO can entrain respiratory rhythms and is anatomically positioned to influence both breathing dynamics and vocal motor output \citep{castellucci2018evidence, smith2013brainstem, feldman2013understanding}. However, its role in generating the diversity of adult USV intonation patterns has not been investigated.

Here we use simultaneous whole-body plethysmography and acoustic recordings to systematically characterize the relationship between breath waveform and pitch patterns across the full murine vocal repertoire. We combine this with electromyography (EMG) of respiratory and laryngeal muscles and optogenetic manipulation of iRO neurons to establish that breath waveform modulation, driven by iRO-mediated diaphragm re-activation during expiration, constitutes a primary intonation mechanism for the majority of USV syllable types.

\section{Methods}

\subsection{Animals and Behavioral Recordings}

Male mice ($n = 6$) were placed in whole-body plethysmography chambers modified with ultrasonic microphones. Fresh female urine was introduced as a social stimulus. Breathing (airflow) and USVs (40--120~kHz) were recorded simultaneously. A total of 1,850 vocalizations were analyzed across all animals.

\subsection{Vocalization Classification}

USVs were classified into 10 syllable types based on pitch contour patterns: Down FM, Flat, Short, Chevron, Step Down, Multi, Step Up, Up FM, Complex, and Two-syllable. Classification used automated spectral analysis with manual verification.

\subsection{Intonation Analysis}

For each vocalization, the instantaneous Pearson correlation between expiratory airflow and pitch was computed. Positive correlation ($r > 0$) indicates positive intonation (pitch tracks airflow); negative correlation ($r < 0$) indicates negative intonation (pitch anticorrelated with airflow). Statistical comparisons used one-way ANOVA with Sidak's post-hoc correction across syllable types.

\subsection{Electromyography}

EMG electrodes were implanted in the diaphragm (inspiratory muscle), thyroarytenoid, and cricothyroid (laryngeal muscles) in a separate cohort. Recordings were performed simultaneously with plethysmography and acoustic monitoring.

\subsection{Optogenetic Experiments}

AAV CreON-FlpON-ChR2::YFP was injected bilaterally into the iRO of Penk-Cre;Vglut2-Flp mice ($n = 5$). Light stimulation at various frequencies (including 10~Hz) was delivered via implanted optical fibers.

\section{Results}

\subsection{Basic Vocalization Parameters}

Table~\ref{tab:basic_params} summarizes fundamental parameters of USV production during courtship-driven vocalization bouts.

\begin{table}[H]
\centering
\caption{Basic vocalization parameters. Values are mean $\pm$ SD across $n = 6$ animals. Paired $t$-test for USV vs.\ non-USV breath comparisons.}
\label{tab:basic_params}
\begin{tabular}{lccc}
\toprule
Parameter & USV Breaths & Non-USV Sniffs & $p$-value \\
\midrule
Instantaneous frequency (Hz) & $7.5 \pm 0.8$ & $9.2 \pm 0.6$ & $0.030$ \\
Peak inspiratory flow (a.u.) & $1.42 \pm 0.18$ & $1.21 \pm 0.14$ & $0.010$ \\
Peak expiratory flow (a.u.) & $1.08 \pm 0.12$ & $1.05 \pm 0.11$ & $0.270$ \\
Inspiratory time (ms) & $52.3 \pm 5.1$ & $50.8 \pm 4.7$ & $0.400$ \\
Expiratory time (ms) & $85.7 \pm 8.3$ & $81.2 \pm 7.6$ & $0.180$ \\
Ti/Te ratio & $0.61 \pm 0.06$ & $0.63 \pm 0.05$ & $0.250$ \\
Syllables per breath & $1.00$ (88\%) & -- & -- \\
\bottomrule
\end{tabular}
\end{table}

USV breaths were slightly slower than neighboring non-USV sniffs ($p = 0.03$) and had larger peak inspiratory flow ($p = 0.01$), but similar timing parameters and expiratory flow.

\subsection{Syllable Type Distribution}

Table~\ref{tab:syllable_dist} shows the distribution of the 10 syllable types across the 1,850 analyzed vocalizations.

\begin{table}[H]
\centering
\caption{USV syllable type distribution ($n = 1{,}850$ total vocalizations across 6 mice). Percentage of total and mean count per mouse $\pm$ SD.}
\label{tab:syllable_dist}
\begin{tabular}{lcc}
\toprule
Syllable Type & \% of Total & Count per Mouse (mean $\pm$ SD) \\
\midrule
Down FM & $27.3$ & $84.2 \pm 18.7$ \\
Flat & $22.1$ & $68.2 \pm 15.3$ \\
Short & $17.4$ & $53.7 \pm 12.8$ \\
Chevron & $10.2$ & $31.5 \pm 8.9$ \\
Step Down & $7.8$ & $24.1 \pm 7.2$ \\
Multi & $5.9$ & $18.2 \pm 5.8$ \\
Step Up & $3.8$ & $11.7 \pm 4.1$ \\
Up FM & $3.1$ & $9.6 \pm 3.5$ \\
Complex & $1.6$ & $4.9 \pm 2.1$ \\
Two-syllable & $0.8$ & $2.5 \pm 1.3$ \\
\bottomrule
\end{tabular}
\end{table}

\subsection{Intonation Mechanism by Syllable Type}

Table~\ref{tab:intonation} shows the mean airflow--pitch correlation coefficient and intonation classification for each syllable type.

\begin{table}[H]
\centering
\caption{Airflow--pitch correlation by USV syllable type. Mean Pearson $r \pm$ SEM. One-way ANOVA with Sidak's post-hoc correction: $^*p < 0.05$, $^{**}p < 0.01$, $^{***}p < 0.001$ (testing $r \neq 0$).}
\label{tab:intonation}
\begin{tabular}{lcccc}
\toprule
Syllable Type & $n$ & Mean $r$ & SEM & Classification \\
\midrule
Down FM & 505 & $0.52$ & $0.08$ & Positive$^{***}$ \\
Flat & 409 & $0.41$ & $0.07$ & Positive$^{***}$ \\
Short & 322 & $0.38$ & $0.09$ & Positive$^{**}$ \\
Chevron & 189 & $0.47$ & $0.10$ & Positive$^{***}$ \\
Multi & 109 & $0.35$ & $0.11$ & Positive$^{*}$ \\
Step Down & 144 & $0.18$ & $0.09$ & Mixed$^{*}$ \\
Step Up & 70 & $-0.39$ & $0.12$ & Negative$^{*}$ \\
Up FM & 57 & $-0.44$ & $0.13$ & Negative$^{*}$ \\
Complex & 30 & $0.05$ & $0.15$ & Combined \\
Two-syllable & 15 & $0.08$ & $0.18$ & Combined \\
\bottomrule
\end{tabular}
\end{table}

Six of 10 syllable types (Down FM, Flat, Short, Chevron, Multi, Step Down) exhibited primarily positive intonation, where pitch tracks expiratory airflow changes. Two types (Step Up, Up FM) showed negative intonation, and two (Complex, Two-syllable) showed mixed patterns.

\subsection{EMG Reveals Diaphragm Re-Activation During Vocalization}

Table~\ref{tab:emg} summarizes the EMG findings during normal breaths versus vocalization breaths.

\begin{table}[H]
\centering
\caption{Muscle activation patterns during normal vs.\ vocalization breaths. Activation index = integrated EMG activity normalized to maximum. Mean $\pm$ SD across recording sessions. Paired $t$-test.}
\label{tab:emg}
\begin{tabular}{llccc}
\toprule
Muscle & Phase & Normal Breaths & USV Breaths & $p$-value \\
\midrule
Diaphragm & Inspiration & $0.82 \pm 0.09$ & $0.85 \pm 0.08$ & $0.430$ \\
Diaphragm & Expiration & $0.08 \pm 0.03$ & $\mathbf{0.54 \pm 0.11}$ & $< 0.001$ \\
Thyroarytenoid & Post-inspiration & $0.71 \pm 0.10$ & $0.68 \pm 0.12$ & $0.520$ \\
Cricothyroid & Post-inspiration & $0.65 \pm 0.11$ & $0.62 \pm 0.13$ & $0.610$ \\
\bottomrule
\end{tabular}
\end{table}

The key finding is a dramatic increase in diaphragm activity during the expiratory phase of vocalization breaths ($0.54 \pm 0.11$ vs.\ $0.08 \pm 0.03$; $p < 0.001$), indicating re-activation of the inspiratory muscle during expiration.

\subsection{iRO Optogenetic Activation Produces Primarily Positive Intonation}

Table~\ref{tab:iro} summarizes the results of iRO optogenetic activation.

\begin{table}[H]
\centering
\caption{iRO optogenetic activation results. Penk$^+$/Vglut2$^+$ neurons in the intermediate reticular tract ($n = 5$ mice). Breathing entrainment assessed at 10~Hz stimulation. Chi-squared test for intonation distribution.}
\label{tab:iro}
\begin{tabular}{lcc}
\toprule
Metric & Evoked USVs & Endogenous USVs \\
\midrule
Syllable types produced & 8 of 10 & 10 of 10 \\
Positive intonation (\%) & $78.3 \pm 6.2$ & $61.5 \pm 5.8$ \\
Negative intonation (\%) & $8.4 \pm 3.1$ & $18.2 \pm 4.3$ \\
Mixed intonation (\%) & $13.3 \pm 4.5$ & $20.3 \pm 5.1$ \\
$\chi^2$ (evoked vs.\ endogenous) & \multicolumn{2}{c}{$p = 0.008$} \\
Breathing entrainment at 10~Hz & Yes & -- \\
\bottomrule
\end{tabular}
\end{table}

iRO activation produced 8 of 10 endogenous syllable types with a significantly higher proportion of positive intonation ($78.3\%$ vs.\ $61.5\%$; $p = 0.008$), consistent with iRO primarily driving breath-mediated pitch modulation. The relative scarcity of negative intonation in evoked USVs suggests that laryngeal-mediated pitch modulation requires an additional neural component.

\subsection{Two-Mechanism Model of Intonation}

Table~\ref{tab:model} summarizes the two-mechanism model.

\begin{table}[H]
\centering
\caption{Two-mechanism model of USV intonation. Each mechanism's contribution estimated from airflow--pitch correlation and EMG/optogenetic data.}
\label{tab:model}
\begin{tabular}{lcccc}
\toprule
Mechanism & Source & Neural Substrate & Syllable Types & Contribution (\%) \\
\midrule
Positive & Breath modulation & iRO $\to$ breathing CPG & 6 of 10 & $61.5 \pm 5.8$ \\
Negative & Laryngeal tension & RAm/laryngeal premotor & 2 of 10 & $18.2 \pm 4.3$ \\
Combined & Both & iRO + laryngeal circuits & 2 of 10 & $20.3 \pm 5.1$ \\
\bottomrule
\end{tabular}
\end{table}

\section{Discussion}

Our results fundamentally revise the understanding of how mice control the pitch of their ultrasonic vocalizations. Rather than relying solely on laryngeal tension modulation, the majority of syllable types use a breath-driven intonation mechanism in which modulation of the expiratory airflow waveform---achieved through re-activation of the inspiratory diaphragm during expiration---directly controls pitch contour. This mechanism is mediated by the iRO brainstem circuit and produces ``positive intonation'' patterns where pitch covaries with airflow.

The diaphragm re-activation we observe during USV production creates a ``mini-breath'' nested within the expiratory phase, modulating airflow velocity and consequently pitch through its effect on subglottal pressure and vocal fold vibration dynamics \citep{titze1988regulation, alipour2002pressure}. This finding resolves the apparent paradox in the literature: while static analyses find weak correlations between breath airflow and pitch, the dynamic relationship within individual vocalizations reveals strong positive correlations for the majority of syllable types.

The two-mechanism model---positive (breath-driven) and negative (laryngeal-driven) intonation---provides a framework for understanding how 10 distinct syllable types can be generated from the combination of two motor control channels. The iRO contributes the breath-shaping component, producing primarily positive intonation, while laryngeal premotor circuits contribute negative intonation through tension modulation independent of airflow \citep{tschida2019vocal, hage2009audio}. Complex and two-syllable types recruit both mechanisms simultaneously, enabling the rich diversity of the murine vocal repertoire.

These findings have implications for studying vocal communication in mouse models of neurodevelopmental disorders, as they identify specific neural circuits and motor mechanisms that can be independently disrupted to produce distinct vocal phenotypes.

\section{Conclusion}

We demonstrate that breath waveform modulation, mediated by iRO-driven diaphragm re-activation during expiration, is the primary intonation mechanism for 6 of 10 mouse USV syllable types. Combined with laryngeal tension modulation, these two mechanisms generate the full diversity of the murine vocal repertoire. This two-mechanism model provides a new framework for understanding vocal motor control in mice and other mammals.

\bibliographystyle{plainnat}
\bibliography{references}

\end{document}
